\chapter{Running and Debugging Make}\label{ch:mk-running}

A Makefile that works is only half the job. The other half is knowing what make
is doing when it does something you did not expect --- which, given that its
entire reasoning is invisible, needs tooling.

\section{Command-Line Options}\label{sec:mk-flags}

\begin{keys}[0.24]
	\cmd{make}             & Build the default goal (\autoref{sec:mk-goal})        \\
	\cmd{make target}      & Build one target                                      \\
	\cmd{make -j8}         & Run up to eight recipes at once                       \\
	\cmd{make -n}          & \vocab{Dry run}: print the recipes, run none          \\
	\cmd{make -B}          & Rebuild everything, ignoring timestamps               \\
	\cmd{make -k}          & Keep going after an error, building what still can be \\
	\cmd{make -C dir}      & Change to \file{dir} first                            \\
	\cmd{make -f build.mk} & Use a file not named \file{Makefile}                  \\
	\cmd{make -s}          & Silent: do not echo recipes                           \\
	\cmd{make -r}          & Drop the built-in rule database                       \\
	\cmd{make -t}          & \vocab{Touch} targets instead of building them        \\
	\cmd{make -q}          & Ask only whether the target is up to date; say so in the exit status \\
\end{keys}

\begin{moral}
	\cmd{make -n} first, whenever a Makefile is unfamiliar or you are about to run
	a target called something like \texttt{deploy}. It prints exactly what would
	run and touches nothing.
\end{moral}

\begin{note}
	\cmd{make -B} is the blunt instrument for ``it should have rebuilt and it did
	not''. It is a diagnostic, not a fix: if \cmd{-B} produces a working binary and
	a plain \cmd{make} does not, you have a missing edge in the graph, and
	\autoref{sec:mk-depgen} is where to look.
\end{note}

\section{Parallel Builds}\label{sec:mk-parallel}

\cmd{make -j} is the reason to have written a correct dependency graph. Make
can run any two recipes simultaneously when neither depends on the other, and on
a modern machine that is most of them.

\begin{shell}
$ time make
real    0m48.2s
$ make clean && time make -j8
real    0m9.7s
\end{shell}

\begin{keys}[0.24]
	\cmd{make -j}          & Unlimited jobs. Rarely what you want --- it will start hundreds \\
	\cmd{make -j8}         & Eight at a time                                        \\
	\cmd{make -j\$(nproc)} & One per core                                           \\
	\cmd{make -l4}         & Start new jobs only while the load average is below 4   \\
	\cmd{make -O}          & Group each job's output, so the log stays readable      \\
\end{keys}

\begin{moral}
	Parallel make is where an incomplete dependency graph finally shows itself. A
	serial build hides a missing edge, because the accidental order happens to be
	right; with \cmd{-j8} the order changes every run, and the build fails
	intermittently. \emph{A Makefile that only works without \cmd{-j} is not
		working.}
\end{moral}

\begin{eg}[The classic race]
\begin{makecode}
all: parser.o lexer.o        # both need parser.h, which is generated

parser.h parser.c: grammar.y
    bison -d -o parser.c $<
\end{makecode}
	Serially, \file{parser.o} is built first and generates the header before
	\file{lexer.o} needs it. In parallel both start at once and one of them fails
	--- sometimes. The fix is to state the edge that was always true:
\begin{makecode}
lexer.o parser.o: parser.h
\end{makecode}
\end{eg}

\begin{note}
	Two recipes writing the same file is the other common parallel bug, and
	\cmd{-j1} will never reveal it. If a build fails only occasionally under
	\cmd{-j}, do not retry until it passes --- the graph is wrong, and the same
	gap will eventually produce a silently stale binary.
\end{note}

\section{Printing from Inside a Makefile}\label{sec:mk-trace}

\begin{keys}[0.30]
	\texttt{\$(info \dots)}    & Print a message while reading the file; expansion continues \\
	\texttt{\$(warning \dots)} & The same, prefixed with the file and line number            \\
	\texttt{\$(error \dots)}   & Print and stop immediately                                  \\
\end{keys}

\begin{makecode}
$(info SRC is [$(SRC)])
$(info OBJ is [$(OBJ)])

ifndef PREFIX
$(error PREFIX is not set -- try `make PREFIX=/usr/local')
endif
\end{makecode}

\begin{note}
	These run at \emph{read} time, not when a recipe runs, so they fire even for
	targets you did not ask for. That is what makes them the right tool for
	inspecting a variable: the square brackets in \texttt{[\$(SRC)]} are there so
	you can see leading and trailing whitespace, which is the usual reason a list
	is not what you expected.
\end{note}

\begin{eg}[A general variable-printing rule]
\begin{makecode}
print-%:
    @echo '$* = [$($*)]'
\end{makecode}
\begin{shell}
$ make print-CFLAGS
CFLAGS = [-Wall -Wextra -std=c17 -g]
$ make print-OBJ
OBJ = [main.o parser.o util.o]
\end{shell}
	Worth keeping in every Makefile you maintain.
\end{eg}

\section{Inspecting the Database}\label{sec:mk-inspect}

\begin{keys}[0.30]
	\cmd{make -p}             & Every variable, rule and built-in, with their origins \\
	\cmd{make -p | less}      & \dots{} which is several thousand lines, so page it   \\
	\cmd{make -d}             & Full decision trace: why each target was or was not rebuilt \\
	\cmd{make --debug=b}      & Basic tracing only --- much more readable than \cmd{-d} \\
	\cmd{make --debug=v}      & Verbose: which files were considered                  \\
	\cmd{make --trace}        & Print each recipe with the line that caused it        \\
	\cmd{make -n -p}          & Inspect without building anything                     \\
\end{keys}

\begin{eg}
	\cmd{make --debug=b} answers the question that matters most often:
\begin{shell}
$ make --debug=b app
Considering target file 'main.o'.
 Prerequisite 'parser.h' is newer than target 'main.o'.
Must remake target 'main.o'.
\end{shell}
	or, in the case you are usually investigating:
\begin{shell}
Considering target file 'main.o'.
 Prerequisite 'main.c' is older than target 'main.o'.
No need to remake target 'main.o'.
\end{shell}
	The second output, when you were expecting a rebuild, says the graph is missing
	an edge rather than that make is broken.
\end{eg}

\begin{note}
	\cmd{make --trace} is the gentler everyday tool: it annotates each recipe with
	the Makefile line and the target that triggered it, which is usually enough to
	find the rule responsible without wading through \cmd{-d}.
\end{note}

\section{Error Messages You Will Meet}\label{sec:mk-errors}

\begin{keys}[0.34]
	\texttt{missing separator}          & A recipe line begins with spaces instead of a tab (\autoref{sec:mk-tab}) \\
	\texttt{No rule to make target 'x'} & \file{x} is needed, does not exist, and no rule matches. Usually a misspelt filename \\
	\texttt{nothing to be done for 'all'} & The target exists and is up to date, but its recipe is empty --- often a missing tab again, or a phony target you forgot to declare \\
	\texttt{'x' is up to date}          & Make is right and you meant a phony target (\autoref{sec:mk-phonyt}) \\
	\texttt{Recursive variable references itself} & \texttt{=} where you wanted \texttt{:=} (\autoref{sec:mk-assign}) \\
	\texttt{commands commence before first target} & A tab-indented line before any rule --- usually a continuation that lost its backslash \\
	\texttt{Circular x <- y dependency dropped} & The graph has a cycle; make breaks it and carries on, so the result is unpredictable \\
	\texttt{Error 1 (ignored)}          & A recipe failed on a line prefixed with \texttt{-} \\
	\texttt{warning: overriding recipe for target 'x'} & Two rules give \texttt{x} a recipe; the later one wins and the earlier is discarded \\
\end{keys}

\begin{note}
	\texttt{make: *** [Makefile:12: app] Error 1} is not a make error at all ---
	it is make reporting that the command on line 12 exited non-zero. The real
	message is above it, from the compiler. Scroll up.
\end{note}

\begin{moral}
	Nearly every confusing make session resolves to one of three questions: is it a
	tab? is the target phony? and is the edge actually in the graph?
	\cmd{cat -A}, \texttt{.PHONY}, and \cmd{make --debug=b} answer them in that
	order.
\end{moral}
